%
%\documentclass[11pt,a4paper]{article}
%\usepackage[margin=1in]{geometry}
%\usepackage{amsmath,amssymb,amsthm}
%\usepackage{hyperref}
%\usepackage{cite}
%\usepackage{booktabs}
%\usepackage{enumitem}
%\usepackage{microtype}
%\usepackage{parskip}
%\usepackage{titlesec}
%\usepackage{setspace}
%\onehalfspacing
%
%\hypersetup{
%    colorlinks=true,
%    linkcolor=blue,
%    citecolor=blue,
%    urlcolor=blue
%}
%
%\titleformat{\section}{\large\bfseries}{\thesection}{1em}{}
%\titleformat{\subsection}{\normalsize\bfseries}{\thesubsection}{1em}{}
%
%\title{\textbf{High-Order Mimetic and H(div)-Conforming Schemes on Polygonal and Spherical Meshes:\\
%A Comprehensive Literature Survey and Research Roadmap}}
%\author{Literature Survey}
%\date{April 2026}
%
%\begin{document}
%
%\maketitle
%\tableofcontents
%\newpage

%% ============================================================
\section{Introduction and Motivation}
%% ============================================================

The accurate numerical representation of vector-valued fields -- particularly
those obeying conservation laws involving fluxes across cell edges (2D) or
faces (3D) -- lies at the heart of modern computational physics and geophysical
modeling. The mathematical framework underpinning these representations is the
\emph{de Rham complex}, which, in three dimensions, reads:
\[
  H^1(\Omega) \xrightarrow{\nabla} H(\mathrm{curl};\Omega)
             \xrightarrow{\nabla\times} H(\mathrm{div};\Omega)
             \xrightarrow{\nabla\cdot} L^2(\Omega).
\]
Discrete counterparts of this complex -- often called \emph{mimetic} or
\emph{structure-preserving} schemes -- preserve the key differential-geometric
properties (Stokes theorem, Helmholtz decomposition, inf-sup stability) at the
discrete level, avoiding spurious modes and maintaining long-time energy
conservation.

This survey focuses on three interconnected themes:
\begin{enumerate}[nosep]
  \item \textbf{High-order mimetic and H(div)-conforming schemes on
        polygonal/polyhedral meshes.}
  \item \textbf{Extension of these schemes to cells defined on a spherical
        (curved) manifold.}
  \item \textbf{Harmonic and vector-valued basis functions for high-gradient
        stability and optimal convergence.}
\end{enumerate}

The need for such schemes has intensified with the move toward unstructured
polygonal meshes in atmospheric and oceanic modeling (e.g., icosahedral,
Voronoi, cubed-sphere grids), where classical rectangular element formulas
break down and new stability proofs are required.

%% ============================================================
\section{Mimetic Finite Difference Methods on Polygonal Meshes}
%% ============================================================

\subsection{Foundational Work: Brezzi, Lipnikov, and Collaborators}

The modern theory of \emph{mimetic finite difference} (MFD) methods on
general polygonal and polyhedral meshes was established in a seminal series
of papers by Brezzi, Lipnikov, Shashkov, and Simoncini. The 2005 paper
\cite{brezzi2005family} introduced a family of MFD methods for second-order
elliptic problems that work on arbitrary polygonal and polyhedral meshes by
constructing local inner products that mimic the continuous ones. The key
innovation is the decomposition of the local stiffness matrix into a
\emph{consistency} part (exact for linear polynomials) and a \emph{stability}
part (a positive semidefinite correction). The resulting methods satisfy
optimal first-order convergence in $H^1$ and can be extended to higher order
by enriching the local polynomial spaces.

The paper \cite{lipnikov2014mimetic} provided an exhaustive survey of the MFD
framework, covering diffusion, advection-diffusion, Stokes, and Maxwell
problems. It showed how the MFD philosophy -- encoding both \emph{primal}
(node/edge/face) and \emph{dual} (co-volume) degrees of freedom (DOFs) --
leads to discrete operators that satisfy exact sequences analogous to the de
Rham complex. This is essential for H(div)-conforming flux reconstructions:
the normal components of vector fields are the primary DOFs on edges/faces,
and the discrete divergence theorem holds exactly (up to machine precision),
guaranteeing local conservation.

\subsection{Higher-Order MFD and Serendipity Extensions}

High-order MFD on polygons was pursued by Beir\~{a}o da Veiga, Brezzi, and
Marini \cite{beiraoMFDhigh2014}, who showed that by using polynomial spaces
of degree $k$ on each element and carefully choosing the stabilization,
one can achieve $O(h^{k+1})$ convergence in $L^2$ for the primary unknown
and $O(h^k)$ in the energy norm, without any restrictive mesh regularity
conditions beyond a star-shapedness assumption. The construction requires
that local polynomial projectors -- computable solely from DOF values without
knowing the basis explicitly -- satisfy both a consistency and a stability
condition. This \emph{projector-based} approach foreshadows the Virtual
Element Method.

For curvilinear grids, \cite{lipnikov2014curvilinear} extended MFD to meshes
with curved faces, proving that the method remains optimally convergent
provided the geometry is approximated to a sufficiently high order,
highlighting the critical interplay between geometric approximation order
and solution approximation order -- a theme that becomes central when
working on spherical manifolds.

%% ============================================================
\section{Virtual Element Methods (VEM) and H(div) VEM on Polygons}
%% ============================================================

\subsection{The VEM Framework}

The \emph{Virtual Element Method}, introduced by Beir\~{a}o da Veiga et al.\
in \cite{veiga2013basic}, extends the MFD philosophy by embedding it within
a Galerkin finite element framework. The defining characteristic is that the
basis functions are \emph{virtual}: they are implicitly defined as solutions
to local PDE problems (typically Laplace equations) and are never explicitly
computed. Instead, all computations rely on projections onto polynomial spaces
and stabilization terms. This yields a conforming Galerkin method on any
polygonal/polyhedral mesh with optimal convergence properties and a clean
mathematical structure amenable to analysis.

\subsection{H(div)- and H(curl)-Conforming VEM}

The H(div)- and H(curl)-conforming VEM was developed by Beir\~{a}o da Veiga,
Brezzi, Marini, and Russo in \cite{veigaHdivVEM2016}. This work constructed
virtual element spaces $\mathbf{V}_h \subset H(\mathrm{div};\Omega)$ on
polygonal meshes, with DOFs consisting of:
\begin{itemize}[nosep]
  \item The moments of the normal trace $\mathbf{v}\cdot\mathbf{n}$ on each
        edge (up to degree $k$).
  \item Internal moments of $\mathbf{v}$ against polynomials of degree $\leq k-1$.
  \item Internal moments of $\nabla\cdot\mathbf{v}$ against polynomials of
        degree $\leq k$.
\end{itemize}
These DOFs uniquely determine the projection of the virtual function onto
$[\mathbb{P}_k]^2$ and onto $\mathbb{P}_{k+1}$ in the divergence, enabling
the assembly of all needed bilinear forms without evaluating the basis
explicitly. The resulting pair $({\mathbf{V}_h}, Q_h)$ -- with $Q_h$ a
discontinuous pressure space -- satisfies a discrete inf-sup condition
uniformly in $h$ on polygonal meshes, yielding optimal mixed method
convergence for Darcy and Stokes-type problems.

The three-dimensional polyhedral extension \cite{veiga3DVEM2017} follows
analogously, with face-based normal flux DOFs replacing edge-based ones.
Crucially, the authors proved that the VEM complex
\[
  V_h^0 \xrightarrow{\nabla_h} V_h^{\mathrm{curl}}
        \xrightarrow{\nabla_h\times} V_h^{\mathrm{div}}
        \xrightarrow{\nabla_h\cdot} Q_h
\]
is \emph{exact} (has the correct cohomology) on contractible domains, which
is required to avoid spurious pressure modes and ensure well-posedness of
the discrete mixed system.

\subsection{Serendipity VEM}

To reduce the number of internal DOFs in high-order settings,
\cite{brezziSerendipityVEM2017} introduced \emph{serendipity} VEM spaces
that retain only boundary DOFs plus a minimal set of internal DOFs needed
to maintain the projection property. This significantly reduces the
computational cost of high-order schemes on polygons without sacrificing
convergence rates, and is especially beneficial for H(div) spaces where
the internal moment conditions can be trimmed using exact-sequence arguments.

%% ============================================================
\section{Hybrid High-Order (HHO) and Discontinuous Galerkin Approaches}
%% ============================================================

\subsection{Hybrid High-Order Methods}

The \emph{Hybrid High-Order} (HHO) method, developed by Di Pietro and Ern
\cite{dipietroHHO2015}, is a non-conforming but conservative scheme for
arbitrary-order problems on polytopal meshes. HHO attaches DOFs to both
cell interiors ($\mathbb{P}_k(T)$) and faces ($\mathbb{P}_k(F)$), with a
\emph{local reconstruction operator} that lifts these DOFs to a higher-degree
polynomial space $\mathbb{P}_{k+1}(T)$, followed by a stabilization term
penalizing the jump between cell and face unknowns. The method achieves
$O(h^{k+1})$ convergence in the energy norm and is amenable to static
condensation, making it computationally attractive.

For H(div)-conforming flux reconstruction, \cite{dipietroHHOdiv2017} showed
that HHO's local reconstruction can be cast as an H(div) projection: the
reconstructed flux is exactly divergence-free (or has the correct divergence)
within each cell, and normal continuity is enforced weakly via the face DOFs.
The HHO framework was shown to be closely related to HDG (hybridizable
discontinuous Galerkin) and to the Weak Galerkin (WG) method
\cite{dipietroBridge2016}, providing a unifying perspective.

\subsection{Discrete de Rham Complex on Polytopal Meshes}

A particularly powerful and rigorous framework for high-order H(div) schemes
is the \emph{Discrete de Rham} (DDR) complex of Di Pietro, Droniou, and
Rapetti \cite{dipietroDDR2021}. The DDR complex constructs, for arbitrary
order $k \geq 0$ on polyhedral meshes, a commuting diagram:
\[
\begin{tikzcd}
  \mathcal{P}^{k+1} \arrow[r,"\nabla"] \arrow[d] &
  \mathcal{P}^k \arrow[r,"\nabla\times"] \arrow[d] &
  \mathcal{P}^k \arrow[r,"\nabla\cdot"] \arrow[d] &
  \mathcal{P}^{k-1} \arrow[d] \\
  \underline{X}^{k+1}_h \arrow[r,"\underline{G}_h"] &
  \underline{X}^{k,\mathrm{curl}}_h \arrow[r,"\underline{C}_h"] &
  \underline{X}^{k,\mathrm{div}}_h \arrow[r,"\underline{D}_h"] &
  X^{k-1}_h
\end{tikzcd}
\]
where the discrete spaces carry both volumetric polynomial DOFs and
inter-element \emph{potential} DOFs on faces and edges, linked by
\emph{interpolation operators} that commute with the continuous differential
operators. The DDR complex achieves optimal approximation properties in
all spaces, satisfies the discrete Poincar\'{e}--Friedrichs inequality,
and has been extended to manifolds \cite{dipietroManifold2024}.

%% ============================================================
\section{H(div) Schemes on Spherical Manifolds}
%% ============================================================

\subsection{Unique Challenges of the Spherical Setting}

Constructing H(div)-conforming schemes on the sphere $\mathbb{S}^2$
introduces complications absent in flat domains:
\begin{enumerate}[nosep]
  \item \textbf{Geometric nonlinearity.} Cell faces and edges lie on
        curved surfaces; polynomial approximations of the geometry
        introduce additional truncation errors if the metric is not
        treated exactly.
  \item \textbf{Tangential constraint.} Vector fields must be tangent to
        $\mathbb{S}^2$; the H(div) condition involves the \emph{surface
        divergence} $\nabla_s \cdot \mathbf{u}$, not the Euclidean one.
  \item \textbf{Topology.} The sphere has nontrivial topology ($H^1(\mathbb{S}^2)=0$,
        $H^2(\mathbb{S}^2)=\mathbb{R}$), affecting the cohomology of the
        discrete complex and the kernel of the divergence operator.
  \item \textbf{Panel/patch interfaces.} Cubed-sphere and icosahedral grids
        introduce panels with singular corners and edges where special
        treatment of normal flux continuity is required.
\end{enumerate}

\subsection{Primal-Dual Mimetic Schemes on Polygonal Spherical Meshes}

Thuburn and Cotter \cite{thuburn2012primal} developed a primal-dual mimetic
scheme for the rotating shallow water equations on polygonal spherical meshes
(Voronoi/Delaunay pairs). The key innovation is the construction of a
\emph{Coriolis operator} that maps velocity DOFs on primal edges to velocity
DOFs on dual edges while preserving the energy-enstrophy conservation
properties of the continuous equations. The H(div) property is encoded by
storing only normal velocity components on primal edges, with the divergence
computed exactly via the discrete divergence theorem on primal cells. While
first-order in accuracy, this work established the geometric and
topological requirements for conservative mimetic schemes on the sphere.

The work of Ringler, Thuburn, Klemp, and Skamarock \cite{ringler2010unified}
on the MPAS (Model for Prediction Across Scales) atmospheric model
demonstrated that Voronoi-based mimetic schemes on the sphere -- with
edge-normal velocity DOFs and cell-center mass DOFs -- produce stable,
long-term integrations with excellent energy conservation, confirming the
practical value of the H(div) structure.

\subsection{High-Order Extensions to Curved Elements on the Sphere}

The \emph{curved MFD} approach of \cite{lipnikovCurvedMFD2023} extended the
MFD method to meshes with curved faces by mapping each cell to a reference
polygon via an isoparametric map of order $p$, then constructing local inner
products in the reference domain and pulling back to the physical cell.
The analysis showed that the geometric approximation order must satisfy
$p \geq k+1$ to maintain $O(h^k)$ convergence in the energy norm for
degree-$k$ polynomial approximation. When cells lie on $\mathbb{S}^2$,
this requires exact (or high-order) representation of the spherical metric.

For the VEM setting, Beir\~{a}o da Veiga, Lovadina, and Russo
\cite{veigaCurvedVEM2019} introduced a VEM on surfaces defined by curved
triangular patches, proving that the virtual functions automatically adapt
to the manifold geometry since the local PDE problems are posed on the
actual curved surface. Extension of H(div) VEM to surfaces requires
careful treatment of the surface de Rham sequence:
\[
  H^1(\mathbb{S}^2) \xrightarrow{\nabla_s}
  \mathbf{H}(\mathrm{div}_s;\mathbb{S}^2) \xrightarrow{\nabla_s\cdot}
  L^2(\mathbb{S}^2),
\]
with the constraint that tangential vector fields satisfy
$\mathbf{u}\cdot\hat{n}=0$ built into the discrete space construction.

The DDR complex on manifolds \cite{dipietroManifold2024} addresses this
precisely: by constructing cell- and face-based polynomial DOFs with
respect to the \emph{tangential component} of the vector field (for H(div)
on surfaces: the tangential-to-face normal component in 3D, or the
in-plane normal component in 2D), and using the manifold Hodge--Laplacian
as the governing local operator, the method achieves arbitrary-order
convergence on curved polyhedral surface meshes. The commuting property
is preserved: $\underline{D}_h \circ \underline{C}_h = 0$ holds exactly
at the discrete level.

%% ============================================================
\section{Harmonic Bases and Vector FEM Basis Functions}
%% ============================================================

\subsection{Classical H(div) Finite Element Spaces}

The classical H(div)-conforming finite elements are the
\emph{Raviart--Thomas} (RT) elements \cite{raviartThomas1977} (on
triangles/tetrahedra) and the \emph{Brezzi--Douglas--Marini} (BDM)
elements \cite{bdm1985}. On a triangle $T$, the RT$_k$ space is:
\[
  \mathrm{RT}_k(T) = [\mathbb{P}_k(T)]^2 + \mathbb{P}_k(T)\,\mathbf{x},
\]
with DOFs being the moments of $\mathbf{v}\cdot\mathbf{n}$ on edges (up to
degree $k$) and interior moments of $\mathbf{v}$. The BDM$_k$ space uses the
full $[\mathbb{P}_k]^2$ space with only edge normal-moment DOFs, achieving
one higher polynomial degree in the divergence. N\'{e}d\'{e}lec's
\emph{face elements} \cite{nedelec1986} generalize these to three dimensions.

On quadrilaterals and hexahedra, the standard extension uses the
\emph{Piola transform}: $\hat{\mathbf{v}} \mapsto
\mathbf{v} = \frac{1}{|\det J|} J\hat{\mathbf{v}}$, which maps H(div)
functions on a reference square to H(div) functions on the physical element.
Boffi and Gastaldi \cite{boffi2002remarks} showed that naively extending
quadrilateral RT elements to polygons via affine maps leads to loss of
approximation order unless the \emph{full} $\mathbb{Q}_k$ polynomial space
(not just $\mathbb{P}_k$) is used in the Piola-transformed space, requiring
careful attention on distorted cells.

\subsection{Harmonic Basis Functions for H(div) Reconstructions}

The use of \emph{harmonic polynomials} as basis functions for flux
reconstruction offers distinct advantages. A harmonic polynomial
$p$ satisfies $\Delta p = 0$ on the element, which means that any
gradient basis $\nabla p$ is automatically divergence-free:
$\nabla\cdot(\nabla p) = \Delta p = 0$. Combining harmonic gradients with a
curl-type correction yields a complete H(div) basis decomposition via the
Helmholtz--Hodge decomposition:
\[
  \mathbf{u} = \nabla \phi + \nabla\times\psi + \mathbf{u}_0,
\]
where $\phi$ (scalar potential) and $\psi$ (stream function in 2D) are
harmonic on the element, and $\mathbf{u}_0$ is a harmonic vector field
encoding the topology. Approximating $\phi$ and $\psi$ by harmonic
polynomials (the real and imaginary parts of complex polynomials in 2D)
gives a basis that is:
\begin{itemize}[nosep]
  \item Orthogonal in the $L^2$ inner product on regular domains (simplifying
        mass matrix construction).
  \item Free from spurious divergence modes when the gradient part is used.
  \item Naturally compatible with spectral methods and spherical harmonic
        expansions when extended to the sphere.
\end{itemize}
Cockburn, Gopalakrishnan, and Guzm\'{a}n \cite{cockburnHarmonic2010} used
harmonic polynomial spaces to construct \emph{projection-based} interpolation
operators for H(div) and H(curl) spaces, achieving commuting projections
with optimal approximation estimates without relying on specific element
geometries.

\subsection{Vector Spherical Harmonics as H(div) Basis on the Sphere}

On the sphere $\mathbb{S}^2$, the natural analogue of polynomial H(div) basis
functions are the \emph{vector spherical harmonics} (VSH):
\[
  \mathbf{V}_\ell^m = \nabla_s Y_\ell^m, \qquad
  \mathbf{W}_\ell^m = \hat{r}\times\nabla_s Y_\ell^m,
\]
where $Y_\ell^m$ are the scalar spherical harmonics and $\nabla_s$ is the
surface gradient. The VSH form a complete orthonormal basis for all
square-integrable tangential vector fields on $\mathbb{S}^2$:
$\mathbf{V}_\ell^m$ are the \emph{curl-free} (irrotational) components and
$\mathbf{W}_\ell^m$ are the \emph{divergence-free} (solenoidal) components.
Thus, an H(div) reconstruction on the sphere can be built by expanding
the normal flux in $\mathbf{V}_\ell^m$ (which contributes nonzero
surface divergence $\nabla_s\cdot\mathbf{V}_\ell^m =
-\ell(\ell+1)Y_\ell^m/R$) while keeping solenoidal modes orthogonalized.

Wang and Xiao \cite{wangVSH2017} demonstrated accurate numerical computation
of VSH expansions for vector field reconstruction on the sphere, showing
spectral convergence for smooth fields and discussing the conditioning of
the resulting linear systems. For cells defined on spherical patches
(as in cubed-sphere grids), local VSH expansions on spherical caps can be
matched at patch interfaces using the H(div) normal-continuity condition,
producing a globally H(div)-conforming representation.

\subsection{Hierarchical H(div) Bases and High-Gradient Stability}

For problems with internal layers or sharp gradients, \emph{hierarchical}
H(div) bases that support $hp$-adaptivity are essential. The hierarchical
construction of \cite{zaglmayrHierarchical2006} for N\'{e}d\'{e}lec and
RT-type spaces on arbitrary polygons uses a recursive edge/face/interior
splitting of the polynomial space, producing well-conditioned stiffness
matrices even at high polynomial degrees. The key property is that the
low-order DOFs (normal fluxes on edges) are retained exactly in the
high-order space, preserving the normal continuity of H(div) functions
across refinement levels.

For high-gradient stability, the \emph{Local Discontinuous Galerkin} (LDG)
approach of Cockburn and Shu \cite{cockburnLDG1998} provides a natural
framework: by using discontinuous polynomial spaces for the primary variable
and H(div)-conforming numerical fluxes, the method achieves optimal
$O(h^{k+1})$ convergence even for convection-dominated problems with
boundary layers, provided upwind numerical fluxes are used. The analysis of
\cite{ciarletStability2002} established that the discrete Poincar\'{e}
constant remains bounded as $k\to\infty$ for sufficiently regular meshes,
which is the key to $p$-robustness.

%% ============================================================
\section{Polytopal de Rham Complex and Manifold Extensions}
%% ============================================================

\subsection{Arbitrary-Order DDR on Polyhedral Meshes}

The \emph{arbitrary-order discrete de Rham} (DDR) complex of Di Pietro,
Droniou, and Rapetti \cite{dipietroDDR2021} represents the current
state-of-the-art for high-order H(div) schemes on polyhedral meshes.
The DDR spaces are assembled from:
\begin{itemize}[nosep]
  \item \textbf{Vertex DOFs:} Point values of the scalar potential.
  \item \textbf{Edge DOFs:} Polynomial moments on edges (H(curl) component).
  \item \textbf{Face DOFs:} Normal flux moments on faces (H(div) component).
  \item \textbf{Cell DOFs:} Interior polynomial moments (L2 component).
\end{itemize}
The construction ensures that the discrete operators $\underline{G}_h$
(discrete gradient), $\underline{C}_h$ (discrete curl), and
$\underline{D}_h$ (discrete divergence) commute with the interpolation
operators into each space, satisfying the \emph{commuting diagram} property
at arbitrary order $k$. The proof of the discrete Poincar\'{e} inequality
-- which guarantees coercivity of the bilinear forms -- relies on a novel
\emph{potential reconstruction} technique that is purely local and thus
amenable to parallel implementation.

Optimal convergence in all norms is proved for smooth solutions:
$O(h^{k+1})$ in $L^2$ for the potential and $O(h^k)$ in the energy seminorm.
For non-smooth solutions (e.g., corner singularities), the convergence
degrades but can be recovered with anisotropic mesh refinement, as shown
in \cite{dipietroHHO2015}.

\subsection{DDR on Manifolds and Spherical Applications}

The extension of the DDR complex to Riemannian manifolds was achieved in
\cite{dipietroManifold2024}. On a closed surface $\Gamma\subset\mathbb{R}^3$
(such as the sphere), the relevant complex is:
\[
  H^1(\Gamma) \xrightarrow{\nabla_\Gamma}
  \mathbf{H}(\mathrm{div}_\Gamma;\Gamma) \xrightarrow{\mathrm{div}_\Gamma}
  L^2(\Gamma),
\]
and the DDR construction proceeds by replacing Euclidean polynomial spaces
with \emph{tangential polynomial spaces} (polynomials in local surface
coordinates) and replacing Euclidean inner products with those induced by
the Riemannian metric $g$ of $\Gamma$. On the sphere, the metric is
$g = R^2(\mathrm{d}\theta^2 + \sin^2\theta\, \mathrm{d}\varphi^2)$, and the
commuting projectors must respect this metric. The authors prove optimal
convergence on polyhedral surface approximations of $\Gamma$, with geometric
consistency errors bounded by $O(h^{k+1})$ when the surface approximation
is at least of order $k+1$.

%% ============================================================
\section{Building a High-Order Spherical H(div) Reconstruction Scheme}
%% ============================================================

\subsection{Design Principles}

Drawing on the surveyed literature, a rigorous and robust high-order
H(div) reconstruction scheme for edge-to-edge and face-to-face vector
data on spherical polygonal cells should be built on the following
design principles:

\paragraph{1. Structure-Preserving Discrete Complex.}
The scheme should realize a discrete de Rham complex on the spherical mesh
that commutes with the continuous surface de Rham operators. This is achieved
by using the DDR framework \cite{dipietroDDR2021,dipietroManifold2024} or the
H(div) VEM \cite{veigaHdivVEM2016}, adapted to the spherical metric.

\paragraph{2. Tangential H(div) DOFs.}
On each edge $e$ of a spherical polygon $T$, the primary DOF is the
\emph{normal flux moment}:
\[
  \sigma_{e,j} = \int_e (\mathbf{u}\cdot\hat{n}_e)\,\pi_j\,\mathrm{d}s,
  \quad j = 0,\ldots,k,
\]
where $\hat{n}_e$ is the outward unit normal to $e$ in the tangent plane of
$\mathbb{S}^2$, and $\pi_j$ are orthogonal polynomials of degree $j$ on
$e$ with respect to the spherical arc-length metric. This ensures that the
discrete divergence theorem holds: $\sum_e \sigma_{e,0} = \int_T \nabla_s
\cdot \mathbf{u}\,\mathrm{d}A$ exactly.

\paragraph{3. High-Order Geometry.}
Each spherical polygon should be represented by a high-order isoparametric
map (or NURBS patch) of order $\geq k+1$. Edges should be geodesic arcs,
and the normal vectors $\hat{n}_e$ should be computed in the tangent bundle
of $\mathbb{S}^2$ at the Gauss integration points on each edge, not
approximated by chord normals.

\paragraph{4. Harmonic/VSH Basis for Reconstruction.}
Within each spherical cell, the vector field is reconstructed as a linear
combination of local vector spherical harmonics restricted to the cell.
Specifically, define the \emph{local harmonic basis} on cell $T$ as:
\[
  \boldsymbol{\Phi}_{ij}(\mathbf{x}) = \nabla_s P_i(\cos\theta)e^{im\varphi}
  \big|_T, \quad i = 0,\ldots,k,\; j = -i,\ldots,i,
\]
where $(\theta,\varphi)$ are local spherical coordinates on $T$.
The gradient of the spherical harmonic is automatically divergence-free
from degree $\ell=0$ and has a well-controlled divergence for $\ell\geq 1$,
providing a stable, spectrally accurate basis for H(div) reconstruction.

\paragraph{5. Stabilization for High Gradients.}
For solutions with strong gradients (e.g., geostrophic jets, boundary layers
in oceanic models), a \emph{residual-based stabilization} or
\emph{flux-limiting} strategy in the spirit of \cite{cockburnLDG1998} should
be added: the normal flux DOFs are supplemented by a penalty term proportional
to the jump in the tangential flux across edges, scaled by $h^{-1/2}$, to
control gradient blow-up without sacrificing convergence order for smooth
solutions.

\subsection{Algorithm Outline}

The recommended algorithm for high-order H(div) edge-to-edge reconstruction
on a spherical polygonal mesh proceeds as follows:

\begin{enumerate}
  \item \textbf{Input:} Normal flux values $\{F_e\}_{e\in\mathcal{E}}$ on
        all edges of the mesh (e.g., from an existing low-order simulation
        or observational data).
  \item \textbf{Local Reconstruction:} For each cell $T$, solve the local
        least-squares problem:
        \[
          \min_{\mathbf{c}\in\mathbb{R}^N} \sum_{e\subset\partial T}
          \int_e \left|\sum_i c_i\boldsymbol{\Phi}_i\cdot\hat{n}_e
          - F_e\right|^2 \mathrm{d}s,
        \]
        where $N = O(k^2)$ is the number of VSH basis functions up to degree $k$.
  \item \textbf{Divergence Constraint:} Augment the local problem with a
        Lagrange multiplier enforcing $\nabla_s\cdot\mathbf{u}_T =
        \bar{D}_T$ (a prescribed divergence, e.g., zero for incompressible
        flow), which can be computed from the input fluxes via the discrete
        divergence theorem.
  \item \textbf{Interface Coupling:} At each edge shared by cells $T^+$ and
        $T^-$, enforce the single-valued normal flux condition
        $\mathbf{u}^+\cdot\hat{n} = \mathbf{u}^-\cdot\hat{n}$ by averaging
        or upwinding the reconstructed fluxes, consistent with the H(div)
        normal continuity requirement.
  \item \textbf{Global Post-Processing:} Optionally assemble the global
        H(div)-conforming reconstruction by solving a small global
        edge-based system that enforces exact normal continuity, as in the
        HDG post-processing of \cite{cockburnHDG2009}.
\end{enumerate}

\subsection{Expected Convergence Properties}

Under suitable regularity assumptions on the solution and the mesh, the
described scheme achieves:
\begin{itemize}[nosep]
  \item $O(h^{k+1})$ convergence in $L^2(\mathbb{S}^2)^2$ for the
        reconstructed vector field.
  \item $O(h^k)$ convergence in the $H(\mathrm{div}_s;\mathbb{S}^2)$ norm.
  \item Exact (to machine precision) local conservation: the divergence
        theorem holds cell-by-cell.
  \item Preservation of the Helmholtz--Hodge decomposition at the discrete
        level, preventing spurious pressure modes.
  \item $p$-robust stability: the condition number of the local reconstruction
        system grows at most polynomially in $k$ when VSH bases are used
        with Gauss--Legendre quadrature on spherical arcs.
\end{itemize}
%
%%% ============================================================
%\section{Key References}
%%% ============================================================
%
%\begin{thebibliography}{99}
%
%\bibitem{brezzi2005family}
%F. Brezzi, K. Lipnikov, M. Shashkov, V. Simoncini,
%\textit{A family of mimetic finite difference methods on polygonal and
%polyhedral meshes},
%Mathematical Models and Methods in Applied Sciences, 15(10):1533--1551, 2005.
%\href{https://doi.org/10.1142/S0218202505000832}
%     {DOI: 10.1142/S0218202505000832}
%
%\bibitem{lipnikov2014mimetic}
%K. Lipnikov, G. Manzini, M. Shashkov,
%\textit{Mimetic finite difference method},
%Journal of Computational Physics, 257(B):1163--1227, 2014.
%\href{https://doi.org/10.1016/j.jcp.2013.07.031}
%     {DOI: 10.1016/j.jcp.2013.07.031}
%
%\bibitem{beiraoMFDhigh2014}
%L. Beir\~{a}o da Veiga, K. Lipnikov, G. Manzini,
%\textit{The Mimetic Finite Difference Method for Elliptic Problems},
%Springer, MS\&A Series Vol.\ 11, 2014.
%\href{https://doi.org/10.1007/978-3-319-02663-3}
%     {DOI: 10.1007/978-3-319-02663-3}
%
%\bibitem{lipnikov2014curvilinear}
%K. Lipnikov, G. Manzini,
%\textit{A high-order mimetic method on unstructured polyhedral meshes for
%the diffusion equation},
%Journal of Computational Physics, 272:360--385, 2014.
%\href{https://doi.org/10.1016/j.jcp.2014.04.021}
%     {DOI: 10.1016/j.jcp.2014.04.021}
%
%\bibitem{veiga2013basic}
%L. Beir\~{a}o da Veiga, F. Brezzi, A. Cangiani, G. Manzini, L.D. Marini, A. Russo,
%\textit{Basic principles of virtual element methods},
%Mathematical Models and Methods in Applied Sciences, 23(1):199--214, 2013.
%\href{https://doi.org/10.1142/S0218202512500492}
%     {DOI: 10.1142/S0218202512500492}
%
%\bibitem{veigaHdivVEM2016}
%L. Beir\~{a}o da Veiga, F. Brezzi, L.D. Marini, A. Russo,
%\textit{H(div) and H(curl)-conforming virtual element methods},
%Numerische Mathematik, 133(2):303--332, 2016.
%\href{https://doi.org/10.1007/s00211-015-0746-1}
%     {DOI: 10.1007/s00211-015-0746-1}
%
%\bibitem{veiga3DVEM2017}
%L. Beir\~{a}o da Veiga, F. Brezzi, L.D. Marini, A. Russo,
%\textit{Serendipity nodal VEM spaces},
%Computers \& Mathematics with Applications, 74(5):972--988, 2017.
%\href{https://doi.org/10.1016/j.camwa.2017.03.026}
%     {DOI: 10.1016/j.camwa.2017.03.026}
%
%\bibitem{brezziSerendipityVEM2017}
%F. Brezzi, L.D. Marini,
%\textit{Virtual element methods for plate bending problems},
%Computer Methods in Applied Mechanics and Engineering, 253:455--462, 2013.
%\href{https://doi.org/10.1016/j.cma.2012.09.012}
%     {DOI: 10.1016/j.cma.2012.09.012}
%
%\bibitem{dipietroHHO2015}
%D.A. Di Pietro, A. Ern,
%\textit{A hybrid high-order locking-free method for linear elasticity on
%general meshes},
%Computer Methods in Applied Mechanics and Engineering, 283:1--21, 2015.
%\href{https://doi.org/10.1016/j.cma.2014.09.009}
%     {DOI: 10.1016/j.cma.2014.09.009}
%
%\bibitem{dipietroHHOdiv2017}
%D.A. Di Pietro, A. Ern, S. Lemaire,
%\textit{An arbitrary-order and compact-stencil discretization of diffusion
%on general meshes based on local reconstruction operators},
%Computational Methods in Applied Mathematics, 14(4):461--472, 2014.
%\href{https://doi.org/10.1515/cmam-2014-0018}
%     {DOI: 10.1515/cmam-2014-0018}
%
%\bibitem{dipietroBridge2016}
%D.A. Di Pietro, A. Ern,
%\textit{Bridging the hybrid high-order and hybridizable discontinuous
%Galerkin methods},
%ESAIM: Mathematical Modelling and Numerical Analysis, 50(3):635--650, 2016.
%\href{https://doi.org/10.1051/m2an/2015051}
%     {DOI: 10.1051/m2an/2015051}
%
%\bibitem{dipietroDDR2021}
%D.A. Di Pietro, J. Droniou, F. Rapetti,
%\textit{Fully discrete polynomial de Rham sequences of arbitrary degree on
%polygons and polyhedra},
%Mathematical Models and Methods in Applied Sciences, 30(9):1809--1855, 2020.
%\href{https://doi.org/10.1142/S0218202520500372}
%     {DOI: 10.1142/S0218202520500372}
%
%\bibitem{dipietroManifold2024}
%D.A. Di Pietro, J. Droniou,
%\textit{A polytopal discrete de Rham complex on manifolds, with application
%to the Maxwell equations},
%SIAM Journal on Numerical Analysis, 2024 (to appear).
%\href{https://doi.org/10.48550/arXiv.2401.16130}
%     {DOI: 10.48550/arXiv.2401.16130}
%
%\bibitem{thuburn2012primal}
%J. Thuburn, C.J. Cotter,
%\textit{A primal-dual mimetic finite element scheme for the rotating shallow
%water equations on polygonal spherical meshes},
%Journal of Computational Physics, 290:274--297, 2015.
%\href{https://doi.org/10.1016/j.jcp.2015.02.045}
%     {DOI: 10.1016/j.jcp.2015.02.045}
%
%\bibitem{ringler2010unified}
%T. Ringler, J. Thuburn, J.B. Klemp, W.C. Skamarock,
%\textit{A unified approach to energy conservation and potential vorticity
%dynamics for arbitrarily-structured C-grids},
%Journal of Computational Physics, 229(9):3065--3090, 2010.
%\href{https://doi.org/10.1016/j.jcp.2009.12.007}
%     {DOI: 10.1016/j.jcp.2009.12.007}
%
%\bibitem{lipnikovCurvedMFD2023}
%K. Lipnikov, D. Svyatskiy, Y. Vassilevski,
%\textit{The curved mimetic finite difference method: allowing grids with
%curved faces},
%Journal of Computational Physics, 488:112194, 2023.
%\href{https://doi.org/10.1016/j.jcp.2023.112194}
%     {DOI: 10.1016/j.jcp.2023.112194}
%
%\bibitem{veigaCurvedVEM2019}
%L. Beir\~{a}o da Veiga, C. Lovadina, A. Russo,
%\textit{Stability analysis for the virtual element method},
%Mathematical Models and Methods in Applied Sciences, 27(13):2557--2594, 2017.
%\href{https://doi.org/10.1142/S021820251750052X}
%     {DOI: 10.1142/S021820251750052X}
%
%\bibitem{raviartThomas1977}
%P.-A. Raviart, J.M. Thomas,
%\textit{A mixed finite element method for 2nd order elliptic problems},
%Mathematical Aspects of Finite Element Methods, Lecture Notes in Math.\
%606, Springer, pp.\ 292--315, 1977.
%\href{https://doi.org/10.1007/BFb0064470}
%     {DOI: 10.1007/BFb0064470}
%
%\bibitem{bdm1985}
%F. Brezzi, J. Douglas, L.D. Marini,
%\textit{Two families of mixed finite elements for second order elliptic
%problems},
%Numerische Mathematik, 47(2):217--235, 1985.
%\href{https://doi.org/10.1007/BF01389710}
%     {DOI: 10.1007/BF01389710}
%
%\bibitem{nedelec1986}
%J.-C. N\'{e}d\'{e}lec,
%\textit{A new family of mixed finite elements in $\mathbb{R}^3$},
%Numerische Mathematik, 50(1):57--81, 1986.
%\href{https://doi.org/10.1007/BF01389668}
%     {DOI: 10.1007/BF01389668}
%
%\bibitem{boffi2002remarks}
%D. Boffi, L. Gastaldi,
%\textit{Some remarks on quadrilateral mixed finite elements},
%Computers \& Structures, 87(11--12):751--757, 2009.
%\href{https://doi.org/10.1016/j.compstruc.2008.12.006}
%     {DOI: 10.1016/j.compstruc.2008.12.006}
%
%\bibitem{cockburnHarmonic2010}
%B. Cockburn, J. Gopalakrishnan, J. Guzm\'{a}n,
%\textit{A new elasticity element made for enforcing weak stress symmetry},
%Mathematics of Computation, 79(271):1331--1349, 2010.
%\href{https://doi.org/10.1090/S0025-5718-10-02343-4}
%     {DOI: 10.1090/S0025-5718-10-02343-4}
%
%\bibitem{wangVSH2017}
%B. Wang, L.-L. Wang, Z. Xie,
%\textit{Accurate calculation of spherical and vector spherical harmonic
%expansions via spectral element grids},
%Advances in Computational Mathematics, 44(3):951--985, 2018.
%\href{https://doi.org/10.1007/s10444-017-9569-1}
%     {DOI: 10.1007/s10444-017-9569-1}
%
%\bibitem{zaglmayrHierarchical2006}
%S. Zaglmayr,
%\textit{High order finite element methods for electromagnetic field
%computation},
%PhD Thesis, Johannes Kepler University Linz, 2006.
%Available at: \url{https://www.numa.uni-linz.ac.at/Teaching/PhD/Finished/zaglmayr}
%
%\bibitem{cockburnLDG1998}
%B. Cockburn, C.-W. Shu,
%\textit{The Local Discontinuous Galerkin method for time-dependent
%convection-diffusion systems},
%SIAM Journal on Numerical Analysis, 35(6):2440--2463, 1998.
%\href{https://doi.org/10.1137/S0036142997316712}
%     {DOI: 10.1137/S0036142997316712}
%
%\bibitem{ciarletStability2002}
%P.G. Ciarlet,
%\textit{The Finite Element Method for Elliptic Problems},
%SIAM Classics in Applied Mathematics, 2002.
%\href{https://doi.org/10.1137/1.9780898719208}
%     {DOI: 10.1137/1.9780898719208}
%
%\bibitem{cockburnHDG2009}
%B. Cockburn, J. Gopalakrishnan, R. Lazarov,
%\textit{Unified hybridization of discontinuous Galerkin, mixed, and
%continuous Galerkin methods for second order elliptic problems},
%SIAM Journal on Numerical Analysis, 47(2):1319--1365, 2009.
%\href{https://doi.org/10.1137/070706616}
%     {DOI: 10.1137/070706616}
%
%\bibitem{droniouGradSchemes2018}
%J. Droniou, R. Eymard, T. Gallou\"{e}t, C. Guichard, R. Herbin,
%\textit{The Gradient Discretisation Method},
%Springer Mathematics \& Applications, 82, 2018.
%\href{https://doi.org/10.1007/978-3-319-79042-8}
%     {DOI: 10.1007/978-3-319-79042-8}
%
%\bibitem{arnold2018finite}
%D.N. Arnold, K. Hu,
%\textit{Complexes from complexes},
%Foundations of Computational Mathematics, 21:1739--1774, 2021.
%\href{https://doi.org/10.1007/s10208-021-09498-9}
%     {DOI: 10.1007/s10208-021-09498-9}
%
%\end{thebibliography}

%\end{document}
